\section*{Data}

To study intrastate conflict patterns in Nigeria we utilize the ACLED data initiated by \citet{raleigh:etal:2010}. This dataset records armed conflict and protest events in over 60 developing countries. ACLED's \textit{battles} data is used to generate our measure of conflict where $y_{ij,t} = 1$ indicates that a conflict occurred when actor $i$ attacked actor $j$ at time $t$ ($y_{ij,t} = 0$ if no conflict occurred).\footnote{In the appendix, we show results when only including battles from ACLED that resulted in at least one battle death. The results when using a fatality threshold and not are similar, thus we focus on the latter for the remainder of this paper.}$^{, }$\footnote{Some scholars have noted that the ACLED conflict data should be understood as defining symmetric events, since who started a particular battle can at times be difficult to entangle. In the appendix, we show that the results presented in the paper are robust to utilizing a directed or undirected formulation.}  We focus only on armed groups that are engaged in battles for at least 5 years during the 2000-2016 period, which results in a total of 37 armed groups.

The groups included in our battle-data analysis are politically violent actors as defined by the ACLED codebook. Such actors include rebels, militias, and ethnic groups.\footnote{For detailed definitions of these actors, see the on-line ACLED Africa data codebook maintained at \url{www.acleddata.com}.} Additionally, we use secondary resources to confirm the existence of each actor in the data. Overall, the data capture several key types of actors: government armed forces (military and police), insurgent groups (Boko Haram), and ethnic militias (including the Yoruba and Fulani Militias). 

The ACLED data also allows us to explore civilians' role in shaping conflict within the Nigerian system.  Building on the research agenda motivated by macro-level studies \citep{gurr:1970, opp:1988, tarrow:1994, tucker:2007,chenoweth:stephan:2011}, we investigate the link between civilian mobilization and violence between armed actors at the local level by creating a count of the number of protests/riots led by civilians against a given actor at time $t-1$. While there is a robust debate over what causes civilian victimization,\footnote{For example, see \citet{valentino:etal:2004, humphreys:weinstein:2006, downes:2006, kalyvas:2006, salehyan:etal:2014, prorok:appel:2014}.} discussion of the consequences of civilian victimization, particularly as they relate to the potential for future conflict has been more limited \citep{hultman:2007, raleigh:2012}. \citet{berman:matanock:2015} argue that when rebels groups [governments] kill civilians, other civilians are more [less] likely to share information with the government: information allows the government to carry out attacks against insurgents, and the lack of information makes insurgent attacks less likely. \citet{condra:shapiro:2012} find support for this relationship in Iraq.\footnote{\citet{lyall:2009}, however, finds a contrary effect and demonstrates that indiscriminate violence against civilians in Chechnya is associated with a lower risk of rebel attacks, possibly because these attacks deprive insurgents of resources. } To account for both these dynamics we create a count of the number of violent actions that an actor committed against civilians at time $t-1$ and nodal covariates to explain both the probability of an actor sending and receiving a conflictual tie. 

Additionally, we attempt to account for spatial effects in modeling the likelihood of conflict between a particular pair of actors. First, we add a measure of how dispersed a rebel group's activity is across the country. Armed groups whose actions are more dispersed across Nigeria are more likely to come into conflictual contact with others. We include a ``Geographic Spread'' variable as both a sender and receiver covariate.\footnote{For example, to model the spread of group $i$ in period $t$, we calculate this by taking all the events involving group $i$ in the 3 years prior to $t$ and estimate the variance.}  Second, a notable literature has developed explicating the mechanisms through which civil conflict can diffuse between countries.\footnote{For example, see \citet{starr:most:1983,salehyan:2009, salehyan:gleditsch:2007,metternich:etal:2015, braithwaite:2010a, beardsley:2011}.} To account for this effect, we use the ACLED dataset to count the number of conflict events occurring within Nigeria's contiguous neighbors. 

We add three other covariates to the model. First, we create a control for whether both actors are a part of the government. There are only two actors within the 37 that we study in the Nigerian system that would fall under this classification, namely, the Police and Military. We include this control to account for the fact that the probability of an interaction between these two is minimal relative to the other potential dyads in this system.\footnote{An alternative approach would be to include sender and receiver level government indicator variables, our results are similar when using that approach.} Second, we include a binary variable that takes on the value of one if the following year is an election year and zero otherwise. This is to account for the fact that elections in Nigeria are rarely waged peacefully in the ballot box, but instead are typically followed or preceded by episodes of conflict. Most importantly, to examine the affect that Boko Haram's entrance has had on the level of conflict in this system we include a binary variable that takes on the value of one after the Boko Haram insurgency has begun (in 2009) and zero beforehand.\footnote{We choose this date both because 2009 is considered the beginning of the Boko Haram insurgency in most media reports, despite their founding in 2002,  because 2009 features the first battle including Boko Haram in the ACLED Dataset.}

\section*{Modeling Approach}

To estimate conflict in this system in a way that explicitly models interdependencies between actors, we rely on a network based approach that combines the social relations regression model (SRRM) -- for details on this model see \citet{li:loken:2002,dorff:minhas:2016} -- and the latent factor model. Together the SRRM provides a set of additive effects that we can use to capture first and second order dependencies, and the latent factor model provides a set of multiplicative effects that we use to model third order dependencies \citep{minhas:etal:2016:arxiv}. This estimator is referred to as the additive and multiplicative effects (AME) model:

\begin{align}
	\begin{aligned}
		y_{ij,t} &= g(\theta_{ij,t}) \\ 
		&\theta_{ij,t} = \bm\beta_{d}^{\top} \mathbf{X}_{ij,t} + \bm\beta_{s}^{\top} \mathbf{X}_{i,t} + \bm\beta_{r}^{\top} \mathbf{X}_{j,t} + e_{ij,t} \\
		&e_{ij,t} = a_{i} + b_{j}  + \epsilon_{ij} + \alpha(\textbf{u}_{i}, \textbf{v}_{j}) \text{  , where } \\
		&\qquad \alpha(\textbf{u}_{i}, \textbf{v}_{j}) = \textbf{u}_{i}^{\top} \textbf{D} \textbf{v}_{j} = \sum_{k \in K} d_{k} u_{ik} v_{jk} \\ 
	\label{eqn:ame}
	\end{aligned}
\end{align}
where $y_{ij,t}$ represents whether or not conflict occurred between actor $i$ (the sender) and actor $j$ (the receiver) at time $t$. To model our binary dependent variable we employ a latent variable representation of a probit regression framework, in which we model a latent variable, $\theta_{ij}$, using a set of time varying dyadic ($\bm\beta_{d}^{\top} \mathbf{X}_{ij,t}$), sender ($\bm\beta_{s}^{\top} \mathbf{X}_{i,t}$), and receiver covariates ($\bm\beta_{r}^{\top} \mathbf{X}_{j,t}$).

To be able to assess the effect that the entrance of actors such as Boko Haram have on the Nigerian conflict network, we extend the AME framework to handle networks in which the composition of actors change over time. Typically, when employing a latent variable framework to longitudinal networks one assumes that the composition of actors is uniform across time \citep{sarkar:moore:2005,ward:etal:2012,sewell:chen:2015}, and if the composition is not truly uniform the authors arbitrarily make it so by choosing a select group of actors to study. This is obviously problematic in general and especially in the case of intrastate conflict as rebel groups often emerge and dissolve -- meaning that in time $t$ we may have actors $\{i,j,k\}$ in the network, and in time $t+1$ actors $\{i,j,k,l\}$ or just a different set of actors $\{i,j,l\}$. Thus we extend the AME framework to handle networks in which the actor composition of a network changes over time. A change in composition can occur in two ways: 1) an actor dissolves after some time, 2) an actor enters the network after some time. \footnote{Theoretically, an actor can also dissolve and then reappear in the network, but this does not occur in the particular context of the Nigerian conflict system between 2000 and 2016.}

To account for these types of changes to a network, we adjust the estimation procedure of AME such that actor observations only contribute to the likelihood of the model once they enter into the network. This procedure is analogous to that adopted by \citet{huisman:snijders:2003} for the stochastic actor oriented model (SAOM). They limit the set of actors who can change their outgoing relations (or have incoming relations changed with) to only those that are part of an active set at time $t$. We implement a process for the AME framework in which composition changes are modeled as exogenous events such that actors are allowed to enter and leave the network at fixed time points. The utility of this approach is that the baseline probability of a tie shifts in accordance with the number of active actors in the network at time $t$ and the time-varying covariates already included in the model.\footnote{A revised version of the AME model is available on Github.}

The $a_{i}$ and $b_{j}$ in Equation~\ref{eqn:ame} represent sender and receiver random effects that we incorporate from the SRRM framework:

\begin{align}
	\begin{aligned}
		\{ (a_{1}, b_{1}), \ldots, (a_{n}, b_{n}) \} &\simiid N(0,\Sigma_{ab}) \\ 
		\{ (\epsilon_{ij}, \epsilon_{ji}) : \; i \neq j\} &\simiid N(0,\Sigma_{\epsilon}), \text{ where } \\
		\Sigma_{ab} = \begin{pmatrix} \sigma_{a}^{2} & \sigma_{ab} \\ \sigma_{ab} & \sigma_{b}^2   \end{pmatrix} \;\;\;\;\; &\Sigma_{\epsilon} = \sigma_{\epsilon}^{2} \begin{pmatrix} 1 & \rho \\ \rho & 1  \end{pmatrix}
	\label{eqn:srm}
	\end{aligned}
\end{align}
The interpretation of these parameters is straightforward. The sender and receiver random effects are modeled jointly from a multivariate normal distribution to account for correlation in how active an actor is in sending and receiving ties. Heterogeneity in the the sender and receiver effects is captured by $\sigma_{a}^{2}$ and $\sigma_{b}^{2}$, respectively, and $\sigma_{ab}$ describes the linear relationship between these two effects (i.e., whether actors who send [receive] a lot of ties also receive [send] a lot of ties). Beyond these first-order dependencies, second-order dependencies are described by $\sigma_{\epsilon}^{2}$ and a within dyad correlation, or reciprocity, parameter $\rho$.\footnote{Note that the latent parameters being estimated are not indexed by $t$, which is to indicate that they represent the average value over time for a particular actor. The interpretation of $a_{i}$, for instance, is the average sender effect across the sample timeframe after accounting for the exogenous covariates.} 

While the additive effects from the SRRM can deal with first (differing levels of activity across actors) and second order interdependencies (reciprocity), the multiplicative effects are used to deal with third order dependencies. Specifically, this multiplicative effect allows us to model homophily -- the tendency of actors with similar characteristics to form strong relationships than those with differing characteristics -- and stochastic equivalence, the possibility that two actors $i$ and $j$ will have similar relationships with every other actor in the network. An AME model accounts for these third order effects using the multiplicative term: $\alpha(\textbf{u}_{i}, \textbf{v}_{j})=\textbf{u}_{i}^{\top} \textbf{D} \textbf{v}_{j}$. This model posits a latent vector of characteristics $\mathbf{u_{i}}$ and $\mathbf{v_{j}}$ for each sender $i$ and receiver $j$. The similarity or dissimilarity of these vectors will then influence the likelihood of activity, and provides a representation of third order interdependencies \citep{minhas:etal:2016:arxiv}.

The representation of third order interdepedencies is accomplished by a process similar to computing the singular value decomposition (SVD) of the observed network. When taking the SVD we factorize our observed conflict network into the product of three matrices: $\textbf{U}, \textbf{D}, \text{ and }, \textbf{V}$.\footnote{The benefit of using an SVD to factorize our observed network is that it is guaranteed to provide the optimal matrix approximation, given $K$, of the observed network.} This provides us with a low-dimensional representation of our original conflict network in terms of third order dependence patterns arising in the network.\footnote{The dimensions of $\textbf{U}$ and $\textbf{V}$ are $n \times K$ and $\textbf{D}$ is a $K \times K$ diagonal matrix.} Specifically, values in $\textbf{U}$ provide a representation of how stochastically equivalent actors are as senders in a network, or more simply put how similar actors are in terms of who they are initiating battles with. For instance, $\hat{\textbf{u}}_{i} \approx \hat{\textbf{u}}_{j}$ would indicate that actor $i$ and $j$ initiate battles with similar third parties. $\textbf{V}$ provide a similar representation but from the perspective of how similar actors are as receivers. The values in $\textbf{D}$, a diagonal matrix, represent levels of homophily in the network.\footnote{Unlike traditional SVD, in the latent factor model, the singular values are not restricted to be positive, thus allowing us to account for both the presence and absence of homophily.} Parameter estimation in the AME takes place by sampling from the posterior distribution of the full conditionals -- further details on the estimation procedure can be found in the Appendix.
